\begin{frame}{자주 묻는 질문 (2/2)}
  \small
  \begin{itemize}
    \item \textbf{Q. 무방향 그래프에 써도 되나요?} --- 된다. 표준
      기법은 무방향 엣지 1개를 \textbf{방향 링크 2개}로 바꿔주는 것.
      이렇게 하면 ``도 순서에 가깝되, \textbf{누가 나를 가리키느냐를
      가중한}'' 순서가 되며, 일반적으로 도 순위와 동일하지는 않다.
    \item \textbf{Q. $p, q$는 그래프마다 튜닝해야 하나요?} ---
      보편적 최적값은 없다. 두 프리셋 $(1,16), (16,1)$ 도 ``역할
      구조 vs 커뮤니티 구조''에 맞춘 \textbf{출발점}에 불과 ---
      몇 가지 설정을 돌려보고 \textbf{downstream 태스크} 성능이나
      시각화에서 기대된 구조가 보이나로 검증한다.
  \end{itemize}
\end{frame}
